%!TEX root = TheoreticalPhysicsC.tex
% Core of content/c03-hilbertraum, c03.s1 to c03.s3. Sources: docs/plan/script-v0.3/sync-ch2.md

\section{Lecture: Hilbert Space}
The \emph{Hilbert space} will play a major role in describing the states of quantum mechanics. To define this space we first need \emph{complex vector spaces}. They are defined in complete analogy with the real case.

Why \emph{complex} and not real vectors appear here, we do not know: nature is simply built that way, and no better alternative has been found so far.

\subsection{Complex vectors}
A \emph{complex vector} $\vb{v}\in\mathbb{C}^n$ is a column of $n$ complex numbers, its \emph{entries}:
\begin{equation}
	\vb{v} = \begin{pmatrix} v_1 \\ \vdots \\ v_n \end{pmatrix}.
\end{equation}
Multiplication by a scalar $a\in\mathbb{C}$ and addition work entry by entry: $a\vb{v}$ has the entries $a v_j$, and $\vb{v}+\vb{w}$ has the entries $v_j+w_j$. The space of all such vectors with fixed $n$, together with these two operations, is called the \emph{complex vector space} $\mathbb{C}^n$; the number $n$ is called its \emph{dimension}. For the vector space these two rules suffice. The dagger operation and the inner product of the next section are extra structure, and only with them does $\mathbb{C}^n$ become a Hilbert space.

The simplest case, $\mathbb{C}^1$, is boring: it is just the complex numbers. The first interesting example is $\mathbb{C}^2$: two-dimensional complex vectors describe the pure states of a qubit.

\begin{example}
	Let
	\begin{equation}
		\vb{v} = \begin{pmatrix} 1+i \\ 1-i \end{pmatrix}, \qquad
		\vb{w} = \begin{pmatrix} 1 \\ -1 \end{pmatrix}.
	\end{equation}
	Computing entry by entry,
	\begin{equation}
		\vb{v}+\vb{w} = \begin{pmatrix} 2+i \\ -i \end{pmatrix}, \qquad
		i\,\vb{v} = \begin{pmatrix} -1+i \\ 1+i \end{pmatrix}.
	\end{equation}
\end{example}

\subsection{The dagger operation and the inner product}
The \emph{dagger operation} turns a column vector into the row vector with the complex conjugate entries. The \emph{inner product} of two vectors $\vb{v},\vb{w}\in\mathbb{C}^n$ is the product of the row $\vb{v}\adj$ with the column $\vb{w}$, a complex number:
\begin{equation}
	\vb{v}\adj = \begin{pmatrix} v_1^* & \cdots & v_n^* \end{pmatrix}, \qquad
	\vb{v}\adj\vb{w} = \sum_{j=1}^{n} v_j^* w_j.
\end{equation}
The dagger always sits on the first vector.

% fig_vl02_kreuz.tex — Die Kreuz-Operation in der komplexen Ebene.
%
% Adaptation of panel (a) of the course figure c03.fig.kreuz
% (content/c03-hilbertraum/figures/kreuz.tex), redrawn with the script's
% TikZ styles (tpc-macros.sty) so that it builds in both masters.
%
% Geometry from the formulas, not freehand: the entries of
%     v = (2+i, -1+2i)  in C^2
% are drawn as the points (Re v_j, Im v_j), their conjugates
%     v_1^* = 2-i,  v_2^* = -1-2i
% as the points (Re v_j, -Im v_j).  Conjugation is the reflection in the
% real axis; the dashed arrows join each entry to its conjugate.
% One unit = 0.62 cm, as in the course figure.
%
% Filled discs: entries of v.  Open circles: entries of v^dagger.  The
% distinction is carried by fill, not by colour, so it survives greyscale.
% Panel (b) of the course figure (the row-times-column product) is left
% out: the script computes the same product as a display equation.

\begin{center}
\begin{tikzpicture}[x=0.62cm, y=0.62cm]
  \draw[axisstyle] (-2.6,0) -- (3.4,0);
  \draw[axisstyle] (0,-2.7) -- (0,2.8);
  \node[lbl, anchor=west] at (3.45,0) {$\operatorname{Re}$};
  \node[lbl, anchor=south] at (0,2.85) {$\operatorname{Im}$};
  \foreach \t in {-2,-1,1,2} {
    \draw[spacecadet, line width=0.4pt] (\t,-0.08) -- (\t,0.08);
    \draw[spacecadet, line width=0.4pt] (-0.08,\t) -- (0.08,\t);
  }
  \node[lbl, anchor=north] at (1,-0.1) {$1$};
  \node[lbl, anchor=east] at (-0.1,1) {$1$};
  % conjugation = reflection in the real axis
  \draw[birrengray, line width=0.5pt, dashed, -{Stealth[length=4pt]}]
    (2,0.78) -- (2,-0.78);
  \draw[birrengray, line width=0.5pt, dashed, -{Stealth[length=4pt]}]
    (-1,1.78) -- (-1,-1.78);
  % entries of v (filled) and of v^dagger (open)
  \fill[cgblue] (2,1) circle (2.6pt);
  \fill[cgblue] (-1,2) circle (2.6pt);
  \draw[cgblue, line width=0.7pt, fill=white] (2,-1) circle (2.6pt);
  \draw[cgblue, line width=0.7pt, fill=white] (-1,-2) circle (2.6pt);
  \node[lbl, anchor=west] at (2.18,1) {$v_1 = 2+i$};
  \node[lbl, anchor=west] at (2.18,-1) {$v_1^* = 2-i$};
  \node[lbl, anchor=east] at (-1.2,2) {$v_2 = -1+2i$};
  \node[lbl, anchor=east] at (-1.2,-2) {$v_2^* = -1-2i$};
\end{tikzpicture}
\end{center}


The figure shows what conjugation does geometrically: it reflects every entry in the real axis, here for the entries $2+i$ and $-1+2i$. One would not have to conjugate; one could simply lay the column on its side. But conjugation pays off as soon as one multiplies a vector with itself. Since $v_j^* v_j = \abs{v_j}^2$, we have
\begin{equation}
	\vb{v}\adj\vb{v} = \sum_{j=1}^{n} \abs{v_j}^2 \geq 0.
\end{equation}
The inner product of a vector with itself is therefore real and never negative. Without conjugation, row times column for the vector with the entries $1$ and $i$ would give the number $1+i^2=0$, although it is not the zero vector.

\begin{example}
	For $\vb{v}$ and $\vb{w}$ from the last example,
	\begin{equation}
		\vb{v}\adj\vb{w} = \begin{pmatrix} 1-i & 1+i \end{pmatrix} \begin{pmatrix} 1 \\ -1 \end{pmatrix} = -2i, \qquad
		\vb{w}\adj\vb{v} = \begin{pmatrix} 1 & -1 \end{pmatrix} \begin{pmatrix} 1+i \\ 1-i \end{pmatrix} = 2i.
	\end{equation}
	So the order matters; the two results are complex conjugates of each other.
\end{example}

\subsection{Infinitely many entries}
At this point one could already define finite-dimensional Hilbert spaces. But infinite-dimensional vector spaces play a very important role in physics. The reason is a physical one: our best models are fields, and fields have infinitely many degrees of freedom. One can picture such vectors as columns with infinitely many entries, or as functions that assign a complex number $f(x)$ to every point $x$ of the real line; functions describe, for instance, the pure states of a spinless particle in one dimension. If one simply sets $n=\infty$ in the formulas, the inner product can diverge; the web course (chapter ``The Hilbert-space toolkit'') shows how to repair this, by admitting only columns with a finite sum of squared moduli (the space $\ell^2$) and functions with a finite integral of the squared modulus.

A complex vector space with an inner product is almost a Hilbert space; one condition is still missing, completeness, and in $\mathbb{C}^n$ it always holds. Before we define the Hilbert space in the next lecture, we need a notation that makes calculations with inner products much easier: Dirac's notation.
